\section{Text, Numbering, and Questions the Draft Must Not Hide}

The Douay--Rheims/Vulgate tradition divides John 6 differently from most modern
editions after verse 50. What modern editions call 6:51---including both the
living-bread declaration and ``the bread that I will give is my flesh''---is
split between Douay--Rheims 6:51 and 6:52. Modern 6:63, on Spirit and flesh, is
Douay--Rheims 6:64. This study uses modern numbering in its exposition and
gives both systems where the difference could mislead a reader.

The registered Greek apparatus also records variants that matter to exposition
without overturning it. At 6:39 it reports differences around the pronoun and
preposition in the resurrection clause. At 6:55 it reports the adverb/adjective
difference just noted. At 6:63 it reports present versus perfect speech
(``words I speak'' / ``have spoken''). At 6:69 the Byzantine confession
``Christ, the Son of the living God'' differs from the shorter ``Holy One of
God'' reported for Nestle--Aland. None changes the chapter's Eucharistic
vocabulary, but 6:69 must not be quoted as textually uniform.

One larger compositional question remains debated: some modern interpreters
have treated 6:51c--58 as a later Eucharistic expansion because its flesh-and-
blood diction seems to sharpen the preceding bread-and-faith discourse. This
study does not adjudicate redaction history. In the received canonical form,
the verses stand in one discourse and the repeated descent, life, eating, and
resurrection language supplies real continuity. Historical hypotheses must not
be presented as excisions already established by the manuscript record.

